\documentclass[11pt]{article}
\usepackage[a4paper,margin=18mm]{geometry}
\usepackage{fontspec}
\setmainfont{Latin Modern Roman}
\usepackage{amsmath,amssymb,amsthm,amscd,graphicx,color}
\usepackage{xurl}
\usepackage[unicode,hidelinks]{hyperref}
\allowdisplaybreaks
\setlength\emergencystretch{3em}
\numberwithin{equation}{section}

\newtheorem{thm}{Theorem}[section]
\newtheorem{cor}[thm]{Corollary}
\newtheorem{lem}[thm]{Lemma}
 { \theoremstyle{definition}
\newtheorem{Remark}[thm]{Remark} }

\def\tr{\operatorname{tr}}
\def\bC{\mathbf{\overline{C}}}
\def\C{\mathbf{C}}
\def\bc{\check b}
\def\pc{\check p}
\def\uc{\check u}
\def\lc{,\allowbreak\ldots,\allowbreak}
\let\ds\displaystyle
\def\Res{\mathop{\mathrm{Res}}\limits}
\def\ut{\tilde u}
\def\yt{\tilde y}
\def\gl{{\mathfrak{gl}}}
\def\Wr{\mathop{\mathrm{Wr\vrule height0pt width.6pt}}\nolimits}


\begin{document}
\begin{center}\Large A second-order Fuchsian equation in form (2.1), with apparent exponents m\_j+1, Fuchs balance and gamma,delta,gamma-alpha-1,delta-alpha-1 outside 0,...,d-1, is obtained from the hypergeometric equation (2.3) by a degree-d polynomial in z*d/dz\end{center}
\noindent\textbf{Stable ID:} 1772\par
\noindent\textbf{Authors:} Alexandre Eremenko; Vitaly Tarasov\par
\noindent\textbf{Original work:} Fuchsian Equations with Three Non-Apparent Singularities\par
\noindent\textbf{Version:} Official SIGMA 14 (2018) article 058, DOI 10.3842/SIGMA.2018.058; journal PDF and native TeX acquired 2026-10-01, exact bytes pinned by SHA256\par
\noindent\textbf{Selected task scope:} Exactly Theorem2.1 for equation2.1: distinct complex a\_j outside\{0,1\}, positive integers m\_j, d=sum m\_j, complex alpha,beta,gamma,delta with gamma+delta=alpha+beta+1+d, all a\_j apparent, and gamma,delta,gamma-alpha-1,delta-alpha-1 not in\{0,...,d-1\}. A degree-d polynomial Q exists such that the solution space is the image under Q(z*d/dz) of the displayed hypergeometric equation2.3; the nonresonance exclusions give a two-dimensional isomorphism and matching monodromy. No arbitrary-exponent apparent-singularity claim or metric/counting corollary is separately included.\par
\noindent\textbf{Difficulty:} H2: Removing arbitrarily many apparent singularities requires an explicit solution-space intertwiner, not only a monodromy comparison. The author constructs it through residue solutions of a bispectral difference equation, Casorati determinants and adjoint factorization. The proof also handles parameter specialization through the algebraic apparency constraints and establishes that the resulting differential operator is injective on the two-dimensional hypergeometric solution space.\par
\noindent\textbf{Editorial boundary note (not author text):} The complete original Theorem 2.1 and full author Sections 4–5 proof are retained with distinct a\_j, positive integral m\_j, apparency, exact Fuchs relation and all four finite nonresonance exclusions. The full new apparency-polynomial/bounded-fiber and deformation support in Section 3 is bundled as uncounted core; the rational-critical-point tableau enumeration is cited established background. Casorati operators, factorizations/adjoints, residue bispectral solutions, degree-d polynomial Q and transformed independent hypergeometric solutions are all supplied. The neighboring counting theorem and the metric application are not separate selected tasks, and no arbitrary-exponent apparent-singularity assertion is added.\par
\noindent\textbf{Source:} \url{https://sigma-journal.com/2018/058/}\quad DOI: \url{https://doi.org/10.3842/SIGMA.2018.058}\par
\noindent\textbf{License and attribution:} Copyright retained by the named authors. Published by SIGMA under CC BY 4.0: \url{https://creativecommons.org/licenses/by/4.0/}. Source: \url{https://sigma-journal.com/about.html}. Adaptation consists of standalone excerpt formatting, cover and numbering restoration; original author language and mathematical text are retained.\par
\noindent\textbf{Editorial symbol notice (not author text):} The residue construction chooses a solution y with nonzero residue-leading coefficients at the finitely many apparent points; the printed proof says let y be a solution before asserting all polynomial degrees m\_j. This is a routine generic choice within the stated two-dimensional solution space, retained as a bounded scope note. Remark2.3 has a z\textasciicircum{}alpha versus z\textasciicircum{}(alpha+1) indicial-power slip; the actual equation2.3 and explicit independent series F1,F2 in Section5 specify the exponents. Author pages are reproduced unchanged.\par
\medskip\hrule\medskip\noindent\textbf{Author text begins below.}\par
\setcounter{section}{1}
\setcounter{subsection}{0}
\setcounter{subsubsection}{0}
\setcounter{thm}{0}
\setcounter{equation}{1}
% Original source lines 87--128
\section{Statement of results}\label{section2}

Let $a_j$ be distinct complex numbers, $a_j\not\in\{0,1\}$ for $1\leq j\leq k$. Let $m_j$ be positive integers, and $\alpha$, $\beta$, $\gamma$, $\delta$ complex numbers. Consider the differential equation
\begin{gather}\label{1a}
y''-\left( \frac\alpha z+\frac\beta{z-1}+\sum_{i=1}^k \frac{m_i}{z-a_i}\right) y'+\frac{\gamma\delta z^k+\cdots{}}{z(z-1)\prod\limits_{i=1}^k(z-a_i)}y = 0,
\end{gather}
where $\gamma+\delta=\alpha+\beta+1+\sum\limits_{j=1}^km_j$ (Fuchs's condition). The Riemann scheme of this equation is
\begin{gather*}
\left(\begin{matrix}
0&1&a_1&\ldots&a_k&\infty\\
0&0&0&\ldots&0&-\gamma\\
\alpha+1&\beta+1&m_1+1&\ldots&m_k+1&-\delta
\end{matrix}\right).
\end{gather*}
We recall that the first line of the Riemann scheme contains the singularities, and the other two the exponents at each singularity~\cite{Mellin}. Set $d=\sum\limits_{j=1}^km_j$. Assume that the singularities at \smash{$a_1\lc a_k$} are all apparent and
\begin{gather}\label{0d}
\gamma, \delta, \gamma-\alpha-1, \delta-\alpha-1 \not\in \{0,1\lc d-1\}.
\end{gather}

\begin{thm}\label{theorem1}
Under these conditions, there exists a polynomial $Q$ of degree $d=\sum\limits_{j=1}^km_j$ such that every solution of \eqref{1a} has the form $Q\big(z\frac{{\rm d}}{{\rm d}z}\big)F(z) $, where $F$ is a solution of the hypergeometric equation
\begin{gather}\label{2a}
F''-\left( \frac\alpha z+\frac{\beta+d}{z-1} \right) F'+\frac{\gamma\delta}{z(z-1)} F = 0.
\end{gather}
In particular, the monodromy group of equation \eqref{1a} coincides with that of the hypergeometric equation~\eqref{2a}.
\end{thm}

\begin{Remark}\label{remark1} Clearly, conditions \eqref{0d} imply that for any non-zero polynomial $p$ of degree at most $d-1$, the functions $p(z)$ and $z^{\alpha+1}p(z)$ are not solutions of the hypergeometric equation~\eqref{2a}.
\end{Remark}

\begin{Remark}\label{remark2} Under conditions \eqref{0d}, for any non-zero polynomial $Q$ of degree at most $d$ and any non-zero solution $F$ of \eqref{2a}, the function $Q\big(z\frac{{\rm d}}{{\rm d}z}\big)F(z)$ is non-zero. Indeed, a solution $F$ of~\eqref{2a} such that $Q\big(z\frac{{\rm d}}{{\rm d}z}\big)F(z)\equiv 0$, has the following form: there are polynomials $p_1$, $p_2$ of degrees at most $d-1$ and constants $A,B$, such that
\begin{alignat*}{3}
& F(z) = p_1(z) + z^\alpha p_2(z)\qquad && \hbox{\rm if $\alpha$ is not an integer},&\\
& F(z) = p_1(z) + z^\alpha p_2(z) (A+B\log z)\qquad && \hbox{\rm if $\alpha$ is a non-negative integer},&\\
& F(z) = z^\alpha p_1(z) + p_2(z) (A+B\log z)\qquad && \hbox{\rm if $\alpha$ is a negative integer}. &
\end{alignat*}
Then in the first case, both $p_1(z)$ and $z^\alpha p_2(z)$ are solutions of equation~\eqref{2a} in contradiction with Remark~\ref{remark1}. In the second case, $z^\alpha p_2(z)$ is a solution of \eqref{2a} if $B\ne 0$, while for $B=0$, $F(z)$ is a polynomial and $ \deg F<d$ since $Q\big(z\frac{{\rm d}}{{\rm d}z}\big)F(z)=0$. Both options contradict Remark~\ref{remark1}. The third case can be worked out similarly.
\end{Remark}

\begin{thm}\label{theorem2}
For every $\alpha$, $\beta$, $\gamma$, $\delta$, $m_1\lc m_k$, $a_1\lc a_k$, there exist at least one and at most $\prod\limits_{j=1}^k(m_j+1)$ equations~\eqref{1a} with apparent singularities at all points $a_j$. For generic $\alpha$, $\beta$, $\gamma$, $\delta$, $a_1\lc a_k$, the equality holds.
\end{thm}

\setcounter{section}{2}
\setcounter{subsection}{0}
\setcounter{subsubsection}{0}
\setcounter{thm}{5}
\setcounter{equation}{5}
% Original source lines 149--443
\section{Proof of Theorem~\ref{theorem2}}\label{section3}

It will be convenient to use a different normalization for our differential equation. By a linear-fractional transformation of the independent variable in (\ref{1a}) we place the singular points in $\C$, so that they become
\begin{gather*}
(z_1,\ldots,z_n,\infty),
\end{gather*}
where $n=k+3$, and infinity is an apparent singularity with exponent difference $1$. Then we multiply the dependent variable on some function to shift the exponents and obtain the Riemann scheme
\begin{gather}\label{rs}
\left(\begin{matrix}z_1&z_2&\ldots&z_n&\infty\vspace{1mm}\\
\ds\frac{1+\alpha_1}{2}&\ds\frac{1+\alpha_2}{2}&\ds\ldots&
\ds\frac{1+\alpha_n}{2}& 0\vspace{1mm}\\
\ds\frac{1-\alpha_1}{2}&\ds\frac{1-\alpha_2}{2}&\ds\ldots&
\ds\frac{1-\alpha_n}{2}&-1\end{matrix}\right),
\end{gather}
where $\alpha_1=\alpha+1$, $\alpha_2=\beta+1$, $\alpha_3=m_1+1\lc \alpha_{n-1}=m_k+1$, $\alpha_n=\gamma-\delta$. This transformation changes neither the exponent differences nor the projective monodromy group, and the new equation has the form
\begin{gather}\label{1b}
w''+\sum_{j=1}^n\left(\frac{1-\alpha_j^2}{4(z-z_j)^2} +\frac{\beta_j}{z-z_j}\right)w=0.
\end{gather}
Regularity of the point $\infty$ implies
\begin{gather}\label{1}
\sum_{j=1}^n\beta_j=0,\\
\label{2}
\sum_{j=1}^n\beta_jz_j=\sum_{j=1}^\infty(1-\alpha_j^2)/4,
\end{gather}
and the condition that the formal solution at $\infty$ contains no logarithm is
\begin{gather}\label{3a}
\sum_{j=1}^n\beta_jz_j^2=\sum_{j=1}^n z_j\big(1-\alpha_j^2\big)/2.
\end{gather}
When $\alpha_j$ is an integer, the singularity at $z_j$ may be apparent. The test for the apparent singularity can be found, for example in~\cite{TL}. Write our equation near $z_j$ as
\begin{gather*}
(z-z_j)^2w''+\left(\sum_{k=0}^\infty x_{k,j}(z-z_j)^k\right)w=0.
\end{gather*}
The singularity at $z_j$ is apparent if and only if
$\alpha_j=\ell_j$ is an integer,
\begin{gather*}
x_{0,j}=\big(1-\ell_j^2\big)/4,
\end{gather*}
and
\begin{gather}\label{app}
Y_{\ell_j}(x_{1,j}\lc x_{\ell_j,j})=0,
\end{gather}
where for any $\ell$, the polynomial $Y_\ell$ is defined by
\begin{gather*}
Y_{\ell}(x_1\lc x_\ell)=
\left|\begin{matrix}x_1&1\cdot(1-\ell)&0&\ldots&0\\
x_2&x_1&2\cdot(2-\ell)&\ldots&0\\
\vdots&\vdots&\vdots&\ddots&\vdots\\
x_{\ell-1}&x_{\ell-2}&x_{\ell-3}&\ldots&(\ell-1)\cdot(-1)\\
x_\ell&x_{\ell-1}&x_{\ell-2}&\ldots&x_1
\end{matrix}\right|.
\end{gather*}
The matrix has $x_1$ on the main diagonal, $x_2$ on the next diagonal below the main one, etc., and the given numbers on the next diagonal above the main one.

Notice that the polynomial $Y_{\ell}$ has the form
\begin{gather}\label{app2}
Y_\ell(x_1\lc x_\ell)=x_1^\ell+Y^*_\ell(x_1\lc x_\ell),\qquad \deg Y^*_\ell< \ell.
\end{gather}
\begin{lem}\label{lem1}
Suppose that all singularities $z_j$ and exponents $\alpha_j$ are fixed, and that $\alpha_3\lc\alpha_{n-1}$ are integers, $\alpha_j=\ell_j$. Then the number of such equations with apparent singularities at $z_3\lc z_{n-1}$ is finite.
\end{lem}

\begin{proof} We take $\beta_3\lc\beta_{n-1}$ as variables, and express $\beta_1$, $\beta_2$, $\beta_n$ in terms of these variables, using \eqref{1}--\eqref{3a}. (The determinant of this $3\times3$ system with respect to $\beta_1$, $\beta_2$, $\beta_n$ is the Vandermonde determinant, so it is not zero.) Then for each $j\in[3,n-1]$ the variables $x_{k,j}$ become linear functions of $\beta_3\lc\beta_{n-1}$. Noticing that $x_{1,j}=\beta_j$, we obtain that in terms of $\beta_j$, $3\leq\beta_j\leq n-1$, our equations \eqref{app}, \eqref{app2} have the form
\begin{gather}\label{equa}
\beta_j^{\ell_j}=S_j(\beta_3\lc\beta_{n-1}),\qquad 3\leq j\leq n-1,
\end{gather}
where $S_j$ are polynomials of degrees at most $\ell_j-1$. Thus
\begin{gather*}
|\beta_j|\leq C\bigl(1+\|\beta\|^{1-1/\ell_j}\bigr)\leq C_1\bigl(1+\|\beta\|^{1-1/\ell}\bigr),\qquad \ell=\max_j\ell_j,
\end{gather*}
where $\|\beta\|^2:=\beta_3^2+\cdots+\beta_{n-1}^2$. Hence, $\|\beta\|\leq C_2 $. As a bounded algebraic set must be finite, this proves the lemma.

So the condition that singularities $z_3\lc z_{n-1}$ of equation \eqref{1b} are apparent is expressed as $n-3$ equations \eqref{equa} with $n-3$ unknowns $\beta_3\lc\beta_{n-1}$. Adding one extra variable $\beta_0$, we make these equations homogeneous:
\begin{gather*}%\label{system}
\beta_j^{\ell_j}=\beta_0 \widetilde S_j(\beta_0,\beta_3\lc\beta_{n-1}),\qquad 3\leq j\leq n-1,
\end{gather*}
with some homogeneous polynomials $\widetilde S$ of degree $\ell_j-1$. Evidently no non-zero solution can have $\beta_0=0$, so Bezout's theorem \cite[Chapter~IV, Section~1]{Sh} implies that the number of solutions
of~\eqref{equa}, counting multiplicity, equals the product of degrees $\ell_3\cdots\ell_{n-1}=\alpha_3\cdots\alpha_{n-1}$.

Now we show that for generic $\alpha_1$, $\alpha_2$, $\alpha_n $, solutions of \eqref{equa} are simple. To do this, we specify $\alpha_1$, $\alpha_2$, $\alpha_n$ to be large integers, and show that in this case, there are {\em at least} $\alpha_3\cdots\alpha_{n-1}$ solutions. It will follow that for generic $\alpha_1$, $\alpha_2$, $\alpha_n $, solutions are of multiplicity~$1$, because the set of all large integers is Zariski dense.

Solutions of \eqref{equa} correspond to equations \eqref{1b} with trivial local projective monodromy around $z_3\lc z_{n-1}$. The number of such equa\-tions~\eqref{1b} is at least the number of equations~\eqref{1b} with trivial projective monodromy around {\em all} singularities, and ratios of solutions of such equations are rational functions with critical points at $z_1\lc z_n$ of multiplicities $\alpha_j-1$ at $z_j$. Rational functions obtained from the same differential equation are {\em equivalent} in the sense that they are obtained from each other by linear-fractional transformations.
\end{proof}

So we prove

\begin{lem}\label{lem2}
Suppose that
\begin{gather}\label{conds}
\alpha_1+\alpha_n\geq (1/2)\sum_{j=1}^n(\alpha_j-1)+2,\qquad \alpha_2\geq\sum_{j=3}^{n-1}(\alpha_j-1)+1.
\end{gather}
Then the number of equivalence classes of rational functions with critical points at generic points~$z_j$ of multiplicities $\alpha_j-1$, for $1\leq j\leq n$ is the product $\alpha_3\cdots\alpha_{n-1}$.
\end{lem}

As a corollary we obtain that trivial local monodromy around $z_4\lc z_n$ implies that the whole projective monodromy group is trivial, assuming that $\alpha_1$, $\alpha_2$, $\alpha_3$ are large enough in comparison with
other~$\alpha_j$.

\begin{proof}[Proof of Lemma~\ref{lem2}]It is known \cite{EGSV} that classes of rational functions with prescribed generic critical points are enumerated by the Young tableaux of shape $2\times (d-1)$, where
\begin{gather*}
d=\frac12 \biggl(\sum_{j=1}^n(\alpha_j-1)+2\biggr)
\end{gather*}
is the degree of the rational function.

The diagrams are filled with $\alpha_j-1$ numbers $j$ by the usual rules: the entries do not decrease in rows (left to right) and strictly increase in columns. It follows from this rule that $1$'s can be only in the first row on the left, while $n$'s in the second row in the right. Condition~\eqref{conds} implies that two entries other than $1$, $n$ cannot stand in the same column: they occupy the place in the first row on the right, and in the second row on the left. There is no other restriction on placement of these entries (from $2$ to $n-1$), except that they must be non-decreasing in their rows. So the tableau is completely determined by specifying the list of entries from $2$ to $n-1$ in the first row. The possible number of $k$'s ($3\leq k\leq n-1$) in the first row is $\alpha_k$, and the $2$'s occupy the remaining places. The second assumption of the lemma ensures that there are always enough of $2$'s to fill the remaining places. So we have $\alpha_3\cdots\alpha_{n-1}$ possibilities.

This proves Lemma~\ref{lem2} and completes the proof of Theorem~\ref{theorem2}.
\end{proof}

\section{Preliminaries on difference equations}\label{section4}

In the proof of Theorem~\ref{theorem1} we use the tools developed in \cite{MTV} which are explained in this section. Our exposition is formally independent of~\cite{MTV}.

Denote by $\tau$ be the shift operator, $(\tau f)(x)=f(x+1)$ acting on meromorphic functions on the real line. For functions $u_1(x)\lc u_k(x)$, define their Casorati determinant (discrete Wronskian) by
\begin{gather*}
\Wr_k[u_1\lc u_k] = \det\big(\tau^{j-1}u_i\big)_{i,j=1}^k.
\end{gather*}
By convention, $\Wr_0=1$.
\begin{lem}\label{L1}
For given functions $u_1(x)\lc u_k(x)$ such that $\Wr_k[u_1\lc u_k]\ne 0$ and a non-zero function $F(x)$, there exists a unique difference operator of the form
\begin{gather}\label{dif}
D_{u_1\lc u_k;F}:=\tau^{k+1}+\sum_{j=1}^k K_j(x)\tau^j+F(x)
\end{gather}
such that $D_{u_1\lc u_k;F}u_i=0$ for every $i=1\lc k$.
\end{lem}

\begin{proof} For given $u_1\lc u_k$ and $F$, consider a system of linear equations
\begin{gather*}
\sum_{j=1}^kK_j \tau^ju_i=-\tau^{k+1}u_i-Fu_i,\qquad i=1\lc k,
\end{gather*}
for the coefficients $K_1\lc K_k$ of $D_{u_1\lc u_k;F}$. The determinant of this linear system equals $\tau\bigl(\Wr_k[u_1\lc u_k]\bigr)\ne 0$, so the system has a unique solution.
\end{proof}

\begin{lem}\label{L2}
For every function $f(x)$, we have
\begin{gather}
D_{u_1\lc u_k; F} f = \frac{\tau\bigl(\Wr_{k+1}[u_1\lc u_k,f]\bigr)+ (-1)^kF\Wr_{k+1}[u_1\lc u_k,f]}{\tau\bigl(\Wr_k[u_1\lc u_k]\bigr)}.\label{e1}
\end{gather}
\end{lem}

\begin{proof} Denote by $D$ the difference operator whose action on $f$ is given by the right-hand side of \eqref{e1}. Then $Du_i=0$ for every $i=1\lc k$, and $D$ is of the form \eqref{dif} with some coefficients $K_1\lc K_k$. Hence, $D=D_{u_1\lc u_k; F}$, by Lemma~\ref{L1}.
\end{proof}

\begin{lem} \label{L3}$ D_{u_1\lc u_k; F}=(\tau-g_{k+1}) (\tau-g_k) \cdots (\tau-g_1)$, where
\begin{gather*}
g_i = \frac{\tau f_i}{f_i},\qquad f_i = \frac{\Wr_i[u_1\lc u_i]}{\Wr_{i-1}[u_1\lc u_{i-1}]},
\qquad i=1\lc k,
\end{gather*}
and
\begin{gather*}
g_{k+1} = \frac{(-1)^{k+1}F\Wr_k[u_1\lc u_k]}
{\tau\bigl(\Wr_k[u_1\lc u_k]\bigr)}.
\end{gather*}
\end{lem}
\begin{proof}Set $F_i=(-1)^i \tau\bigl(\Wr_i[u_1\lc u_i]\bigr)\big/ \Wr_i[u_1\lc u_i]$ for $i=1\lc k$, and $F_{k+1}=F$. Let $D_i=(\tau-g_i) \dots (\tau-g_1)$. We will prove by induction on $i$, that $D_i=D_{u_1\lc u_{i-1}; F_i}$ for any $i=1\lc k+1$.

The base of induction for $i=1$ is immediate. Clearly, $D_i=\tau^i+\sum\limits_{j=1}^{i-1} K_{i,j} \tau^j+F_i$ for some coefficients $K_{i,1}\lc K_{i,i-1}$. Thus by Lemma~\ref{L1}, it remains to show that $D_iu_j=0$ for any $j=1\lc i-1$. Since $D_i=(\tau-g_i) D_{i-1}$, it suffices to show that $D_{i-1}u_j=0$ for any $j=1\lc i-1$. This follows from Lemma~\ref{L2} since $D_{i-1}=D_{u_1\lc u_{i-2}; F_{i-1}}$ by the induction assumption.
\end{proof}

\begin{lem}\label{L4} Let $u_1(x)\lc u_k(x)$ be solutions of the difference equation
\begin{gather}
\sum_{i=0}^{k+1} A_i(x) u(x+i) = 0,\qquad A_0\ne 0,\qquad A_{k+1}\ne 0.\label{e2}
\end{gather}
Let $w(x)$ be a non-zero solution of the difference equation $A_0(x) w(x)=(-1)^{k+1}A_{k+1}(x) w(x+1)$. Then the function
\begin{gather*}
v(x) = \frac{\Wr_k[u_1\lc u_k](x+1)}{A_0(x) w(x)}%\label{e3}
\end{gather*}
satisfies the difference equation
\begin{gather}
\sum_{i=0}^{k+1} A_i(x-i) v(x-i) = 0.\label{e4}
\end{gather}
\end{lem}

\begin{proof} If $\Wr_k[u_1\lc u_k]=0$, the claim is trivial. If $\Wr_k[u_1\lc u_k]\ne 0$, consider the difference operator $D=D_{u_1\lc u_k; K_0}$ of the form \eqref{dif} with $K_i=A_i/A_{k+1}$, $i=0\lc k$. Then the difference equation \eqref{e2} reads $A_{k+1} Du=0$.

Define the difference operator $D^\dag=\tau^{-k-1}+\sum\limits_{i=0}^kK_i(x-i) \tau^{-i}$. Then the difference equa\-tion~\eqref{e4} reads $D^\dag(A_{k+1}v)=0$.

The operator $D^\dag$ is obtained from $D$ by applying the formal conjugation antiautomorphism that sends $\tau$ to $\tau^{-1}$ and the multiplication by $x$ to itself. Thus Lemma~\ref{L3} yields
\begin{gather*}
D^\dag= \big(\tau^{-1}-g_1\big) \cdots \big(\tau^{-1}-g_k\big) \big(\tau^{-1}-g_{k+1}\big),
\end{gather*}
where one takes $F=A_0/A_{k+1}$. Straightforwardly, $(\tau^{-1}-g_{k+1}) (A_{k+1}v)=0$, which proves Lemma~\ref{L4}.
 \end{proof}

Now we consider our equation \eqref{1a}. Define the polynomial $A$ and numbers $b_{i,2}$ by
\begin{gather}\label{A}
A(z)=(z-1)\prod_{i=1}^k(z-a_i)= z^{k+1}+\sum_{i=1}^k b_{i,2}z^i+(-1)^{k+1}a_1\cdots a_k,
\end{gather}
and
\begin{gather*}
W(z)=z^\alpha(z-1)^\beta\prod_{i=1}^k(z-a_i)^{m_i}.
\end{gather*}
Multiply equation \eqref{1a} by $z^2A(z)$, so that it becomes
\begin{gather}
A(z) \left(z\frac{{\rm d}}{{\rm d}z}\right)^{ 2}y-B(z) z\frac{{\rm d}y}{{\rm d}z}+C(z)y=0,\label{3}
\end{gather}
where
\begin{gather}%\label{B1}
B(z) = A(z)\bigl(1+zW'(z)/W(z)\bigr)\nonumber\\
\hphantom{B(z)}{} = (\gamma+\delta)z^{k+1}-\sum_{i=1}^k b_{i,1}z^i+(-1)^{k+1}(\alpha+1) a_1\cdots a_k,\label{B2}
\end{gather}
with some numbers $b_{i,1}$, and $C(z)$ is a polynomial of degree $k+1$ such that $C(0)=0$.

Let $x$ be a new variable and $\tau$ be the shift: $(\tau u)(x)=u(x+1)$. The difference equation
\begin{gather}
\left(x^2A(\tau)-xB(\tau)+C(\tau)\right) u(x)=0\label{4}
\end{gather}
is called the {\it bispectral dual} of equation \eqref{3}. Bispectral dual equations are discussed in Theorems~4.1 and~4.2 of \cite{MTV}.

Write \eqref{4} in the form
\begin{gather}
(-1)^{k+1}a_1\cdots a_k x(x-\alpha-1) u(x) +\sum_{i=1}^k b_i(x) u(x+i) \nonumber\\
\qquad{} +(x-\gamma)(x-\delta) u(x+k+1) = 0, \label{5}
\end{gather}
where $b_1\lc b_k$ are quadratic polynomials
\begin{gather*}%\label{bs}
b_i(x)=b_{i,2}x^2+b_{i,1}x+b_{i,0},
\end{gather*}
the constants $b_{i,1}$ and $b_{i,2}$ being defined in \eqref{B2} and~\eqref{A}, respectively.

The equation
\begin{gather}
\big(A\big(\tau^{-1}\big) x^2-B\big(\tau^{-1}\big) x+C\big(\tau^{-1}\big)\big)v(x)=0\label{6}
\end{gather}
is called the {\it formal conjugate} to equation \eqref{4}. Equivalently, \eqref{6} can be written as
\begin{gather}
(-1)^{k+1}a_1\cdots a_k x(x-\alpha-1) v(x) +\sum_{i=1}^k b_i(x-i) v(x-i)\nonumber\\
\qquad{} + (x-k-1-\gamma)(x-k-1-\delta) v(x-k-1) = 0. \label{7}
\end{gather}
If $u_1\lc u_k$ are solutions of equation \eqref{5}, then by Lemma~\ref{L4}
\begin{gather}
v(x) = (a_1\cdots a_k)^{-x-1} \frac{\Gamma(x-\gamma) \Gamma(x-\delta)}{\Gamma(x+1) \Gamma(x-\alpha)}\det\bigl(u_i(x+j)\bigr)_{i,j=1}^k\label{8}
\end{gather}
is a solution of equation \eqref{7}.

Equations \eqref{7} and \eqref{1a} are related as follows: if $v(x)$ solves \eqref{7} and the series
\begin{gather}
f(z) = \sum_{n=-\infty}^\infty v(n+\xi) z^{n+\xi}\label{8a}
\end{gather}
converges for some $\xi $, then $f(z)$ is a solution of equation \eqref{1a}.

\section{Proof of Theorem~\ref{theorem1}}\label{section5}

Let $y(z)$ be a solution of equation \eqref{1a}. Define functions
\begin{gather}
 u_i(x) = \oint_{a_i}\frac{y(z) z^{x-\alpha-2} {\rm d}z}
{(z-1)^{\beta+1}\prod\limits_{j=1}^k(z-a_j)^{m_j+1}}
= 2\pi\sqrt{-1} \Res_{z=a_i} \frac{y(z) z^{x-\alpha-2}}
{(z-1)^{\beta+1}\prod\limits_{j=1}^k(z-a_j)^{m_j+1}},\label{9}
\end{gather}
where the integral is over a small circle around $a_i$ counterclockwise. Clearly, $u_i(x)=a_i^xp_i(x)$, where $p_i(x)$ is a polynomial of degree $m_i$. Since $y(z)$ solves equation~\eqref{3}, the function
\begin{gather}\label{yt}
\yt(z) =
\frac{y(z)}{z^{\alpha+1}(z-1)^{\beta+1}\prod\limits_{j=1}^k(z-a_j)^{m_j+1}} =\frac{y(z)}{zA(z)W(z)}
\end{gather}
satisfies the differential equation
\begin{gather}
\left(z\frac{{\rm d}}{{\rm d}z}\right)^{ 2} \bigl( A(z) \yt\bigr)+ z\frac{{\rm d}}{{\rm d}z}\bigl(B(z) \yt)+C(z) \yt = 0.\label{3t}
\end{gather}
By formulas \eqref{9}, \eqref{yt},{\samepage
\begin{gather*}%\label{uiyt}
u_i(x)=\oint_{a_i}\yt(z) z^{x-1} {\rm d}z,
\end{gather*}
and equation \eqref{3t} yields that $u_1\lc u_k$ are solutions of equation \eqref{4}.}

Use functions \eqref{9} in formula \eqref{8} for a solution $v(x)$ of equation~\eqref{7}. Then
\begin{gather}
v(x) = \frac{\Gamma(x-\gamma) \Gamma(x-\delta)}{\Gamma(x+1) \Gamma(x-\alpha)} Q(x),\label{10}
\end{gather}
where
\begin{gather*}%\label{Q}
Q(x) = (a_1\cdots a_k)^{-x-1}\det\bigl(u_i(x+j)\bigr)_{i,j=1}^k=
 \det\bigl(a_i^{j-1}p_i(x+j)\bigr)_{i,j=1}^k
\end{gather*}
is a non-zero polynomial of degree $d=m_1+\dots+m_k$.

To prove Theorem~\ref{theorem1}, it suffices to show that for any solution $F(z)$ of \eqref{2a}, the function $Q\big(z\frac{{\rm d}}{{\rm d}z}\big)F(z)$ solves \eqref{1a} because under conditions \eqref{0d} the function $Q\big(z\frac{{\rm d}}{{\rm d}z}\big)F(z)$ is non-zero, see Remark~\ref{remark2} after Theorem~\ref{theorem1}. Moreover, it is enough to prove this statement for generic values of $\alpha$, $\gamma$, $\delta$. Indeed, by the proof of Theorem~\ref{theorem2}, equation~\eqref{1a} can be deformed to generic values of $\alpha$, $\beta$, $\gamma$, $\delta$ without moving the singularities at $a_1\lc a_k$ and keeping all of them apparent. To make $F$ and $Q$ change continuously under the deformation, one can parametrize~$F$ and a~solution~$y(z)$ of~\eqref{1a} defining~$Q$ by their data at a fixed regular point.

For generic $\alpha$, $\gamma$, $\delta$, the functions
\begin{gather*}
F_1(z) = \sum_{n=0}^\infty \frac{\Gamma(n-\gamma) \Gamma(n-\delta)}{\Gamma(n+1) \Gamma(n-\alpha)} z^n = \frac{\Gamma(-\alpha)}{\Gamma(-\gamma) \Gamma(-\delta)} \,{}_2F_1(-\gamma,-\delta;-\alpha;z)
\end{gather*}
and
\begin{gather*}
F_2(z)= \sum_{n=0}^\infty
\frac{\Gamma(n+1+\alpha-\gamma) \Gamma(n+1+\alpha-\delta)} {\Gamma(n+2+\alpha) \Gamma(n+1)} z^{n+1+\alpha}\\
\hphantom{F_2(z)}{} = \frac{\Gamma(2+\alpha)} {\Gamma(1+\alpha-\gamma) \Gamma(1+\alpha-\delta)} z^{1+\alpha} \, {}_2F_1(1+\alpha-\gamma,1+\alpha-\delta;2+\alpha;z)
\end{gather*}
are independent solutions of the hypergeometric equation \eqref{2a}, see for example \cite{Mellin}. The functions
\begin{gather*}
f_1(z) = Q\left({z\frac{{\rm d}}{{\rm d}z}}\right) F_1(z) = \sum_{n=0}^\infty \frac{\Gamma(n-\gamma) \Gamma(n-\delta)}{\Gamma(n+1) \Gamma(n-\alpha)} Q(n) z^n,\\
f_2(z) = Q\left({z\frac{{\rm d}}{{\rm d}z}}\right) F_2(z) = \sum_{n=0}^\infty\frac{\Gamma(n+1+\alpha-\gamma) \Gamma(n+1+\alpha-\delta)}{\Gamma(n+2+\alpha) \Gamma(n+1)} Q(n+1+\alpha) z^{n+1+\alpha},
\end{gather*}
coincide with the series~\eqref{8a} for function (\ref{10}) and $\xi=0$ or $\xi=1+\alpha$. Since the series converge for $|z|<1$, both $f_1(z)$ and~$f_2(z)$ are solutions of~(\ref{1a}).

This proves Theorem~\ref{theorem1}.

\begin{thebibliography}{99}
\bibitem[2]{TL}
Cui G., Gao Y., Rugh H.H., Tan L., Rational maps as {S}chwarzian primitives,
 \href{https://doi.org/10.1007/s11425-016-5140-7}{\textit{Sci. China Math.}} \textbf{59} (2016), 1267--1284,
 \href{https://arxiv.org/abs/1511.04246}{arXiv:1511.04246}.

\bibitem[5]{EGSV}
Eremenko A., Gabrielov A., Shapiro M., Vainshtein A., Rational functions and
 real {S}chubert calculus, \href{https://doi.org/10.1090/S0002-9939-05-08048-2}{\textit{Proc. Amer. Math. Soc.}} \textbf{134}
 (2006), 949--957, \href{https://arxiv.org/abs/math.AG/0407408}{math.AG/0407408}.

\bibitem[10]{Mellin}
Iwasaki K., Kimura H., Shimomura S., Yoshida M., From {G}auss to {P}ainlev\'e.
 A modern theory of special functions, \href{https://doi.org/10.1007/978-3-322-90163-7}{\textit{Aspects of Mathematics}}, Vol.~E16, Friedr. Vieweg \& Sohn, Braunschweig, 1991.

\bibitem[13]{MTV}
Mukhin E., Tarasov V., Varchenko A., Bispectral and
 {$({\mathfrak{gl}}_N,{\mathfrak{gl}}_M)$} dualities, discrete versus
 differential, \href{https://doi.org/10.1016/j.aim.2007.11.022}{\textit{Adv. Math.}} \textbf{218} (2008), 216--265,
 \href{https://arxiv.org/abs/math.QA/0605172}{math.QA/0605172}.

\bibitem[17]{Sh}
Shafarevich I.R., Basic algebraic geometry. I.~Varieties in projective space,
 2nd~ed., \href{https://doi.org/10.1007/978-3-642-57908-0}{Springer-Verlag}, Berlin, 1994.

\end{thebibliography}
\end{document}
